\documentclass[11pt]{article}
\usepackage[a4paper,margin=1in]{geometry}
\usepackage{fontspec}
\usepackage[boldfont=false]{xeCJK}
\setCJKmainfont{FandolSong-Regular.otf}
\usepackage{amsmath}
\usepackage{amssymb}
\usepackage{array}
\usepackage{booktabs}
\usepackage{graphicx}
\usepackage{adjustbox}
\usepackage{enumitem}
\setlist[itemize]{topsep=2pt,partopsep=0pt,parsep=0pt,itemsep=2pt,leftmargin=1.4em}
\setlist[enumerate]{topsep=2pt,partopsep=0pt,parsep=0pt,itemsep=2pt,leftmargin=1.6em}
\usepackage{parskip}
\usepackage{fvextra}
\fvset{breaklines=true,breakanywhere=true,fontsize=\small}
\setlength{\parindent}{0pt}

\begin{document}

\begin{center}
  {\LARGE\bfseries Mensuration}\\[0.35em]
  {\large Igcse Mathematics}
\end{center}
\vspace{0.6em}

This handout covers Topic 5, Mensuration (measuring length, area and volume). The Core and Extended content here is almost the same. In the exam, some formulas are given in the List of formulas, but you should still learn them all.

\section*{Units of measure}

We use \textbf{metric} 公制 units for \textbf{mass} 质量 (g, kg), \textbf{length} 长度 (mm, cm, m, km), \textbf{area} 面积, \textbf{volume} 体积 and \textbf{capacity} 容量 (ml, litres — the space inside a container).

To change units, be careful with squares and cubes:

\begin{itemize}
\item Length: $1\text{ m} = 100\text{ cm}$.

\item Area: $1\text{ m}^{2} = 100^{2} = 10\,000\text{ cm}^{2}$.

\item Volume: $1\text{ m}^{3} = 100^{3} = 1\,000\,000\text{ cm}^{3}$.

\item Capacity: $1\text{ litre} = 1000\text{ cm}^{3}$, so $1\text{ m}^{3} = 1000$ litres.

\end{itemize}

\textbf{Worked example.} Convert $3\text{ m}^{2}$ to $\text{cm}^{2}$.

\[
3 \times 10\,000 = 30\,000\text{ cm}^{2}.
\]

\section*{Perimeter and area of basic shapes}

The \textbf{perimeter} 周长 is the distance all the way round a shape. The \textbf{area} is the amount of flat space inside it. Here $b$ is the \textbf{base} 底 and $h$ is the perpendicular \textbf{height} 高.

\begin{center}
\adjustbox{max width=\linewidth}{%
\begin{tabular}{ll}
\toprule
\textbf{Shape} & \textbf{Area} \\
\midrule
\textbf{rectangle} 矩形 & $\text{length} \times \text{width}$ \\
\textbf{triangle} 三角形 & $\tfrac{1}{2} \times b \times h$ \\
\textbf{parallelogram} 平行四边形 & $b \times h$ \\
\textbf{trapezium} 梯形 & $\tfrac{1}{2}(a + b)h$, where $a$ and $b$ are the two parallel sides \\
\bottomrule
\end{tabular}}
\end{center}

\textbf{Worked example.} A trapezium has parallel sides $6\text{ cm}$ and $10\text{ cm}$, and height $4\text{ cm}$. Find its area.

\[
\tfrac{1}{2}(6 + 10) \times 4 = \tfrac{1}{2} \times 16 \times 4 = 32\text{ cm}^{2}.
\]

\section*{Circles}

For a \textbf{circle} 圆 with \textbf{radius} 半径 $r$ (and \textbf{diameter} 直径 $d = 2r$):

\[
\text{circumference} = 2\pi r = \pi d, \qquad \text{area} = \pi r^{2}.
\]

The \textbf{circumference} 圆周 is the distance round the circle.

\textbf{Worked example.} A circle has radius $7\text{ cm}$. Find its circumference and area (leave $\pi$ in the answer).

\[
\text{circumference} = 2\pi \times 7 = 14\pi\text{ cm}, \qquad \text{area} = \pi \times 7^{2} = 49\pi\text{ cm}^{2}.
\]

\section*{Arcs and sectors}

An \textbf{arc} 弧 is part of the circumference. A \textbf{sector} 扇形 is a "pizza slice" between two radii. If the sector angle is $\theta$, the arc and sector are that fraction $\dfrac{\theta}{360}$ of the whole circle:

\[
\text{arc length} = \frac{\theta}{360} \times 2\pi r, \qquad \text{sector area} = \frac{\theta}{360} \times \pi r^{2}.
\]

A small slice is a \textbf{minor sector} 小扇形; the large rest is a \textbf{major sector} 大扇形.

\textbf{Worked example.} Find the \textbf{arc length} 弧长 and area of a sector with angle $90^{\circ}$ and radius $8\text{ cm}$.

The fraction is $\dfrac{90}{360} = \dfrac{1}{4}$, so

\[
\text{arc} = \tfrac{1}{4} \times 2\pi \times 8 = 4\pi\text{ cm}, \qquad \text{area} = \tfrac{1}{4} \times \pi \times 8^{2} = 16\pi\text{ cm}^{2}.
\]

\section*{Surface area and volume of solids}

The \textbf{surface area} 表面积 is the total area of all the outside faces. The volume is the space inside. For these solids ($r$ = radius, $h$ = height):

\begin{center}
\adjustbox{max width=\linewidth}{%
\begin{tabular}{lll}
\toprule
\textbf{Solid} & \textbf{Volume} & \textbf{Surface area} \\
\midrule
\textbf{cuboid} 长方体 & $\text{length} \times \text{width} \times \text{height}$ & add the six faces \\
\textbf{prism} 棱柱 & (\textbf{cross-section} 横截面 area) $\times$ length & — \\
\textbf{cylinder} 圆柱 & $\pi r^{2} h$ & $2\pi r h$ (\textbf{curved surface area} 侧面积) $+\, 2\pi r^{2}$ \\
\textbf{pyramid} 棱锥 & $\tfrac{1}{3} \times \text{base area} \times h$ & — \\
\textbf{cone} 圆锥 & $\tfrac{1}{3}\pi r^{2} h$ & $\pi r l$ (curved) $+\, \pi r^{2}$, where $l$ is the \textbf{slant height} 斜高 \\
\textbf{sphere} 球 & $\tfrac{4}{3}\pi r^{3}$ & $4\pi r^{2}$ \\
\bottomrule
\end{tabular}}
\end{center}

\textbf{Worked example.} A cylinder has radius $5\text{ cm}$ and height $10\text{ cm}$. Find its volume and total surface area (in terms of $\pi$).

\[
\text{volume} = \pi \times 5^{2} \times 10 = 250\pi\text{ cm}^{3}.
\]

\[
\text{surface area} = 2\pi(5)(10) + 2\pi(5)^{2} = 100\pi + 50\pi = 150\pi\text{ cm}^{2}.
\]

\section*{Compound shapes and parts of shapes}

A \textbf{compound shape} 组合图形 is made by joining or cutting simple shapes. Split it into parts you know, then add or subtract.

For "parts" of a circle or solid, take the right fraction. For example, a \textbf{hemisphere} 半球 (half a sphere) has volume

\[
\tfrac{1}{2} \times \tfrac{4}{3}\pi r^{3} = \tfrac{2}{3}\pi r^{3}.
\]

A \textbf{frustum} 平截头体 is a cone or pyramid with its top cut off; find its volume by subtracting the small top cone from the whole cone.

\textbf{Worked example.} Find the area of a shape made of a rectangle $8\text{ cm} \times 5\text{ cm}$ with a semicircle of diameter $5\text{ cm}$ on one end.

The semicircle has radius $2.5\text{ cm}$:

\[
\text{area} = 8 \times 5 + \tfrac{1}{2}\pi (2.5)^{2} = 40 + 3.125\pi \approx 49.8\text{ cm}^{2}.
\]


\end{document}
